% !TeX TS-program = xelatex
% 虚构演示简历 — 仅用于 hellojobs 测试数据
\documentclass{resume-photo}
\usepackage{xeCJK}
\setCJKmainfont{Noto Serif CJK SC}
\usepackage{fontspec}
\usepackage{enumitem}
\setlist[itemize]{itemsep=0pt, topsep=1pt, parsep=0pt, partopsep=0pt}

\ResumeName{刘洋}

\begin{document}

\ResumeContacts{
  138-0004-0004,%
  \ResumeUrl{mailto:liuyang@mail.ustc.edu.cn}{liuyang@mail.ustc.edu.cn},%
  \textnormal{中国科学技术大学 | 计算机应用技术 · 硕士 | 2023—2026}%
}

\ResumeTitle

\noindent 计算机视觉与多模态，熟悉 PyTorch、DETR 系列与模型蒸馏；在工业缺陷检测数据集上 mAP@0.5 提升 6.2 个百分点。注重可观测性与文档沉淀，习惯在评审阶段拆解风险；参与线上复盘与跨团队联调，英语 CET-6，可阅读官方文档与 RFC（演示数据）。

\section{教育背景}
\ResumeItem
{中国科学技术大学}
[\textnormal{计算机学院 | 计算机应用技术 | 硕士}]
[2023.09—2026.06(预计)]

\begin{itemize}
  \item 实验室：媒体计算与智能感知。
\end{itemize}


\ResumeItem
{学术交流与在线课程}
[\textnormal{暑校 / MOOC / 厂商认证（演示）}]
[2022—2025]

\begin{itemize}
  \item 完成《深入理解计算机系统》配套实验与 MIT 6.828 / 6.S081 公开课部分章节跟做。
  \item Coursera / edX 分布式系统、云计算相关专项课程结业证书 4 门（虚构编号）。
  \item 通过 AWS SAA / 阿里云 ACP / Redis 开发者等认证中的 1--2 项（简历测试数据）。
  \item 参加 CCF CCSP / 华为 ICT / 百度之星等学科竞赛集训营并完成全部作业。
\end{itemize}



\section{专业技能}
\begin{itemize}
  \item 掌握 Python、PyTorch，熟悉目标检测、实例分割与半监督学习。
  \item 了解 TensorRT 推理加速与 ONNX 导出流程。
  \item 熟悉 OpenCV、Albumentations 数据增强管线。
  \item 熟悉 Linux 日常运维与 Shell，具备日志检索与 CPU/内存突增初步定位能力。
  \item 掌握 Git / CI 与 Code Review，了解 Docker 与 Kubernetes 基本编排与滚动发布。
  \item 了解 OAuth2 / JWT、接口幂等与 REST 规范，具备单元测试与集成测试习惯（JUnit/pytest 等）。
  \item 能阅读英文技术文档，具备跨团队沟通与需求澄清能力（测试简历）。
\end{itemize}

\section{实习经历}
\ResumeItem
{商汤科技}
[\textnormal{计算机视觉算法实习生}]
[2024.06—2024.12]

\begin{itemize}
  \item \textbf{职责}: 产线 AOI 小样本缺陷检测模型迭代。
  \item \textbf{成果}: 在 500 张标注数据上 Recall@90\%Precision 提升 11\%。
\end{itemize}


\ResumeItem
{百度 | 移动生态}
[\textnormal{后端研发实习生}]
[2023.06—2023.11]

\begin{itemize}
  \item \textbf{检索}: 参与搜索相关性特征工程与离线任务调度稳定性治理。
  \item \textbf{存储}: 优化 HBase Get 批量模式，P95 延迟下降 19\%。
  \item \textbf{平台}: 熟悉内部 BVC 构建与发布流水线。
  \item \textbf{算法对接}: 与算法同学联调 A/B 实验埋点与分流一致性校验。
  \item \textbf{文档}: 维护接口字典与错误码说明，降低联调沟通成本。
\end{itemize}



\section{项目经历}
\ResumeItem
{轻量化检测头设计}
[\textnormal{毕业设计课题}]
[2024.03—至今]

\begin{itemize}
  \item 引入可分离卷积与特征金字塔重加权，参数量下降 40\%，速度提升 22 FPS。
\end{itemize}


\ResumeItem
{低代码表单引擎}
[\textnormal{拖拽 + JSON Schema（演示）}]
[2023.07—2023.12]

\begin{itemize}
  \item \textbf{前端}: 基于 React 实现组件注册表与撤销重做栈。
  \item \textbf{后端}: 动态表单元数据存储与版本迁移脚本。
  \item \textbf{生态}: 导出 JSON 与后端校验规则一致，减少联调返工。
  \item \textbf{采用}: 在内 3 个业务线试点，表单搭建时间缩短 70\%（演示）。
\end{itemize}



\section{个人荣誉}
\begin{itemize}
  \item VALSE 2024 学生志愿者
  \item 中科大 学业奖学金 一等
  \item 全国计算机等级考试四级（网络工程师）（演示）
  \item 校级「优秀学生干部 / 优秀团员」或技术社团骨干（综合测评前 15\%）
  \item 开源贡献：向知名项目提交被合并的 PR（文档/测试/feature）（演示）
\end{itemize}

\end{document}
